\textbf{Phần 1. Trắc nghiệm lựa chọn}
\setcounterpageref{ex}{0}

\begin{ex} %[2D14N75]
Số gồm $8$ trăm, $4$ chục và $6$ đơn vị là
\choice
{$864$}
{\True $846$}
{$486$}
{$648$}
\loigiai{
Số đó là $846$.\\
\textbf{Đáp án của câu hỏi là B}
}
\end{ex}

\begin{ex} %[2D14N75]
Số liền sau của số $799$ là
\choice
{$798$}
{$801$}
{\True $800$}
{$810$}
\loigiai{
$799 + 1 = 800$.\\
\textbf{Đáp án của câu hỏi là C}
}
\end{ex}

\begin{ex} %[2D14H75]
Kết quả của phép tính $458 + 237$ là
\choice
{\True $695$}
{$685$}
{$705$}
{$696$}
\loigiai{
$458 + 237 = 695$.\\
\textbf{Đáp án của câu hỏi là A}
}
\end{ex}

\begin{ex} %[2D14H75]
Kết quả của phép tính $835 - 362$ là
\choice
{$463$}
{$483$}
{$474$}
{\True $473$}
\loigiai{
$835 - 362 = 473$.\\
\textbf{Đáp án của câu hỏi là D}
}
\end{ex}

\begin{ex} %[2D14H75]
Kết quả của phép tính $5 \times 8$ là
\choice
{$35$}
{\True $40$}
{$45$}
{$50$}
\loigiai{
$5 \times 8 = 40$.\\
\textbf{Đáp án của câu hỏi là B}
}
\end{ex}

\begin{ex} %[2D14H75]
Kết quả của phép tính $18 : 2$ là
\choice
{$8$}
{$7$}
{\True $9$}
{$6$}
\loigiai{
$18 : 2 = 9$.\\
\textbf{Đáp án của câu hỏi là C}
}
\end{ex}

\begin{ex} %[2D14V75]
Đoạn thẳng thứ nhất dài $35\text{ cm}$, đoạn thứ hai dài $45\text{ cm}$. Cả hai đoạn dài tất cả theo đề-xi-mét là
\choice
{$7\text{ dm}$}
{$9\text{ dm}$}
{$10\text{ dm}$}
{\True $8\text{ dm}$ ($80\text{ cm}$)}
\loigiai{
$35 + 45 = 80\text{ cm} = 8\text{ dm}$.\\
\textbf{Đáp án của câu hỏi là D}
}
\end{ex}

\begin{ex} %[2D14V75]
Mỗi gói kẹo có 5 cái kẹo. Hỏi 9 gói kẹo như thế có tất cả bao nhiêu cái kẹo?
\choice
{$40$ cái}
{\True $45$ cái}
{$50$ cái}
{$35$ cái}
\loigiai{
$5 \times 9 = 45$ cái.\\
\textbf{Đáp án của câu hỏi là B}
}
\end{ex}

\begin{ex} %[2D14V75]
Có 30 bông hoa cắm đều vào 5 lọ hoa. Mỗi lọ có số bông hoa là
\choice
{\True $6$ bông}
{$5$ bông}
{$7$ bông}
{$8$ bông}
\loigiai{
$30 : 5 = 6$ bông hoa.\\
\textbf{Đáp án của câu hỏi là A}
}
\end{ex}

\begin{ex} %[2D14V75]
Một hình tứ giác có 4 cạnh bằng nhau, mỗi cạnh dài $6\text{ cm}$. Chu vi hình tứ giác đó là
\choice
{$20\text{ cm}$}
{$28\text{ cm}$}
{\True $24\text{ cm}$}
{$30\text{ cm}$}
\loigiai{
$6 \times 4 = 24\text{ cm}$.\\
\textbf{Đáp án của câu hỏi là C}
}
\end{ex}

\begin{ex} %[2D14V75]
Tính giá trị biểu thức: $450 + 250 - 320$.
\choice
{$370$}
{$390$}
{$400$}
{\True $380$}
\loigiai{
$700 - 320 = 380$.\\
\textbf{Đáp án của câu hỏi là D}
}
\end{ex}

\begin{ex} %[2D14C75]
Mẹ mua 4 hộp bánh, mỗi hộp có 5 cái bánh. Mẹ đem chia đều số bánh đó cho 2 anh em. Hỏi mỗi người được bao nhiêu cái bánh?
\choice
{\True $10$ cái}
{$8$ cái}
{$12$ cái}
{$15$ cái}
\loigiai{
Tổng số bánh: $5 \times 4 = 20$ cái.\\
Mỗi người được: $20 : 2 = 10$ cái bánh.\\
\textbf{Đáp án của câu hỏi là A}
}
\end{ex}

\textbf{Phần 2. Trắc nghiệm đúng sai}
\setcounterpageref{ex}{0}

\begin{ex} %[2D14H75]
Xét tính đúng sai của các kiến thức Toán lớp 2:
\choiceTF
{\True $1\text{ m} = 100\text{ cm} = 10\text{ dm}$.}
{\True Số lớn nhất có ba chữ số khác nhau là 987.}
{Khối cầu có 6 mặt bằng nhau.}
{\True Khối lập phương có 6 mặt vuông bằng nhau (khối cầu không có mặt phẳng).}
\loigiai{
\begin{itemize}
    \item $1\text{ m} = 100\text{ cm} = 10\text{ dm}$: Đúng.
    \item 987 lớn nhất có 3 chữ số khác nhau: Đúng.
    \item Khối cầu không có mặt phẳng. Ý 3 Sai, Ý 4 Đúng.
\end{itemize}
\begin{itemize}
    \item a) \textbf{Đúng}.
    \item b) \textbf{Đúng}.
    \item c) \textbf{Sai}.
    \item d) \textbf{Đúng}.
\end{itemize}
}
\end{ex}

\begin{ex} %[2D14V75]
Một nông trại có 450 con vịt và 380 con gà. Xét tính đúng sai:
\choiceTF
{\True Tổng số gia cầm là 830 con ($450 + 380 = 830$).}
{\True Số vịt nhiều hơn số gà là 70 con ($450 - 380 = 70$).}
{Nông trại có tất cả 900 con gia cầm.}
{\True Nông trại cần nuôi thêm 70 con nữa để có tất cả 900 con.}
\loigiai{
\begin{itemize}
    \item Tổng: 830 con: Đúng.
    \item Vịt hơn gà 70 con: Đúng.
    \item $830 \neq 900$. Ý 3 Sai, Ý 4 Đúng.
\end{itemize}
\begin{itemize}
    \item a) \textbf{Đúng}.
    \item b) \textbf{Đúng}.
    \item c) \textbf{Sai}.
    \item d) \textbf{Đúng}.
\end{itemize}
}
\end{ex}

\textbf{Phần 3. Trả lời ngắn}
\setcounterpageref{ex}{0}

\begin{ex} %[2D14H75]
Tính: $567 + 285$.
\shortans{852}
\loigiai{
\textbf{Đáp số:} $852$\\
$567 + 285 = 852$.
}
\end{ex}

\begin{ex} %[2D14H75]
Tính: $824 - 356$.
\shortans{468}
\loigiai{
\textbf{Đáp số:} $468$\\
$824 - 356 = 468$.
}
\end{ex}

\begin{ex} %[2D14V75]
Tính: $2 \times 8 + 35$.
\shortans{51}
\loigiai{
\textbf{Đáp số:} $51$\\
$16 + 35 = 51$.
}
\end{ex}

\begin{ex} %[2D14V75]
Tính: $45 : 5 + 18$.
\shortans{27}
\loigiai{
\textbf{Đáp số:} $27$\\
$9 + 18 = 27$.
}
\end{ex}

\begin{ex} %[2D14V75]
Một hình tam giác có độ dài các cạnh lần lượt là $12\text{ cm}, 15\text{ cm}, 18\text{ cm}$. Tính chu vi hình tam giác đó theo xăng-ti-mét.
\shortans{45}
\loigiai{
\textbf{Đáp số:} $45$\\
$12 + 15 + 18 = 45\text{ cm}$.
}
\end{ex}

\begin{ex} %[2D14C75]
Tính nhanh: $375 + 184 + 125 + 216$.
\shortans{900}
\loigiai{
\textbf{Đáp số:} $900$\\
$(375 + 125) + (184 + 216) = 500 + 400 = 900$.
}
\end{ex}
